\section{Limitations and Threats to Validity}\label{sec:limitations}
We state these at length because several of them bound the conclusions above rather than
merely qualifying them.

\paragraph{One head-training domain}
Every fitted head in this study is fitted on \texttt{market1501/train}. That is what makes the
in-domain versus transfer contrast clean, and it is also the study's largest structural limit:
``the student's gain is concentrated where the probe had labels'' is established with one
source domain. Whether the same asymmetry appears with a head fitted on MSMT17 or on a vehicle
set is not tested here, and it is the first experiment we would add.

\paragraph{MSMT17 is measured for the teachers only}
\Cref{tab:coverage} shows the gap. MSMT17 is the large-scale person transfer target and the
dataset best placed to settle the retention question of \cref{tab:retention} on a
gallery-comparable footing, and the C-RADIOv4 rows for it are not yet run. The retention ratios
we do report are not gallery-matched --- CUHK03-NP's gallery is roughly a third of
Market-1501's --- and should not be placed beside published retention figures without that
caveat.

\paragraph{DINOv3 is a stand-in}
C-RADIOv4's actual DINOv3 teacher is the 7B model, which does not fit on the GPU this
study ran on. The DINOv3 rows bracket the student's capacity from above and below and represent
the teacher \emph{family}; SigLIP2-g is the genuine article. Conclusions about the DINOv3 side
of the bracket are correspondingly weaker than those about the SigLIP2 side.

\paragraph{The CCVID anomaly}
CLIP ViT-B/16 at 87M scores \clipCcvid{} mAP on CCVID, ahead of every larger encoder
here. We have verified that it is not an adapter fault --- the protocol and manifest digests are
shared with every other row in that column --- and we have no mechanism for it. CCVID is the
one dataset in the study whose features pass through a tracklet pooling step, which is the
obvious place to look and which we have not yet ruled in or out. Until it is explained, we do
not draw any conclusion from the CCVID column, and readers should not either.

\paragraph{Two protocols are non-standard, for reasons outside our control}
Occluded-REID ships no standard split and labels no cameras, so its protocol uses the whole set
and cannot apply the same-camera exclusion that every Market-1501 number depends on. VRAI
withholds its test identities, so the only offline-scorable protocol runs over its training
split. That is sound for every row in this paper --- no head here has seen a VRAI image --- and
it would be unsound for any model fine-tuned on VRAI, which is a constraint on who may cite
these rows rather than on the rows themselves.

\paragraph{Absolute values are not comparable across datasets}
Gallery sizes here span roughly a thousand to over thirty thousand images. Every claim we make
is within a column, where all rows share a protocol digest and a manifest digest. The tables
print query and gallery sizes (\cref{tab:datasets}) so that any reader tempted to compare
across columns can see the reason not to.

\paragraph{Throughput measures a machine}
The figures in \cref{tab:cost} come from one thermally throttled consumer GPU and are useful
only as ratios within that table, as the caption's example shows.

\paragraph{Licence heterogeneity}
Of the checkpoints here, only the two C-RADIOv4 weights carry a verified licence permitting
commercial use; the rest are unverified or research-only, as are most of the datasets. Our
substrate reports this per run. It does not change any measurement and it changes the practical
recommendation considerably.

\paragraph{No backbone was trained}
This is a design choice, not an oversight, and it bounds the claims accordingly. Nothing here
speaks to what these encoders do when fine-tuned, and \cref{sec:readout}'s reordering result
should not be assumed to survive fine-tuning --- probe rankings are known not to predict
fine-tuned rankings in general, and testing that here would have cost more compute than the
entire study.
